% Reference entries, T (AN-11).  Law 1: os/iword/core/edit.c iw_ed_tab
% ([Tab] C1/02: to the next stop, every 0.5" relative to the left margin;
% WP recomputes the stop when the text before the tab changes, iw_ed_fixup;
% deleted without a question), spec/iword/wp5_format.h "tabs, indents and
% hard pages", face_a.c (Tab; Shift-, Ctrl- and Alt-Tab act as Tab).
% Shots: ../wordperfect51-decomp mint MINT09, MINT10, MINT21, MINT23.
% GATES: test-pubs-iword test-iword-codes test-iword-wp5-roundtrip test-iword-face test-flair-iword-block
\Gates{test-pubs-iword test-iword-codes test-iword-wp5-roundtrip test-iword-face test-flair-iword-block}
\begin{Feature}{Tab}
Tab moves the text that follows to the next tab stop.  Use it to indent the
first line of a paragraph and to line up short items in columns, such as
the entries of a form.

To tab,
\begin{WPSteps}
\item Press \K{Tab}.  The cursor moves to the next tab stop.
\item Type the text.
\end{WPSteps}
\Documents{wpkey}{Tab}

\BeginReplay{ref-tab}
For example, type \TypeRun{Name:}, press \PressKey{Tab}, type
\TypeRun{Peter Gibbons} and press \PressKey{Enter}; type \TypeRun{Dept:},
press \PressKey{Tab} twice\Keys{Tab} and type \TypeRun{Software}.  Press
\Press{Reveal Codes}{Alt-F3} to see the tab codes:
\WordScreen[rows=14]{ref-tab}{Tabs on the editing screen and in Reveal Codes}
\EndReplay

\Notes
\Topic{Code}
Each Tab puts a \Code{[Tab]} code in the document.  Delete it to remove the
tab; InitechWord deletes it without asking.

\Topic{Shift-Tab, Ctrl-Tab and Alt-Tab}
These act as Tab does.

\Topic{Tab Stops}
The tab stops are every half inch, measured from the left margin: under the
standard margins, every five columns of the editing screen.  You cannot set
other tab stops.  When you change the left margin the stops move with it.
Reveal Codes shows the stops on its ruler.

\Topic{Text Before a Tab}
If you add or delete text before a tab on the same line, the tab moves the
text after it to the next stop after the text, so tabbed columns stay lined
up as long as the text before them fits.

\SeeAlso{Indent; Margins, Left and Right; Reveal Codes}
\end{Feature}

\begin{Feature}{Template}
The function-key template lies above the function keys of your keyboard and
names the feature each key starts.  Each key has four lines on the template:
Ctrl and the key in red, Alt and the key in blue, Shift and the key in
green, the key alone in black.  A line is blank where the combination does
nothing.  The key's name is printed under its lines, and the keys are in
groups of four, as on the keyboard, with a legend of the colors between the
groups.

The panel at the left of the template is a quick list of the keys you use
most that are not function keys: the hard page, hard return, tab, Cancel,
and the Home combinations for the top and end of the document, the start and
end of a line, and the place before the codes.

\Notes
\Topic{Ordering a Template}
Ask your Initech representative for part number 1062-0000-20.

\SeeAlso{Cursor Movement; Status Line}
\end{Feature}
